\section{Hoja de ruta hacia la AGI: del pie de montaña de ChatGPT a la cumbre principal de la acción}

Las secciones anteriores trataron, respectivamente, el techo de los modelos, los entornos de acción y la diferenciación entre empresas. Esta sección entra en la hoja de ruta de Guangmi (广密) como «fundamentalista de la AGI». Él sostiene que el aumento de la inteligencia es la única línea principal y que la inteligencia misma es la mayor aplicación. ChatGPT es solo el pie de la montaña; detrás vienen varias cimas más: Coding, Coding Agent, General Agent, AI for Science, Robotics, entre otras.

\begin{figure}[H]
\centering
\includegraphics[width=0.96\textwidth]{agi-roadmap-mountain.png}
\caption{La cordillera de la ruta hacia la AGI: ChatGPT es solo el pie de la montaña; las cumbres principales que siguen entran, una tras otra, en entornos de acción más exigentes. Diagrama conceptual de elaboración propia, a partir de la conversación de 00:30:18--00:43:00.}
\end{figure}

\begin{knowledgebox}{Lectura de la figura: la cumbre principal pasa del diálogo a la acción}
De ChatGPT a Coding, Coding Agent, General Agent, AI4Science y Robotics, la ruta se acerca cada vez más a la acción real y a la retroalimentación real. La conversación demuestra que el modelo puede expresar inteligencia; los Agent y la ciencia demuestran que el modelo puede usar la inteligencia para modificar su entorno.
\end{knowledgebox}

\subsection{La inteligencia misma es la mayor aplicación}

Esta subsección explica la idea de que «la inteligencia misma es la mayor aplicación». La frase no niega los productos concretos; subraya que todos los productos aprovechan el excedente de capacidad de los modelos. Productos como ChatGPT, Cursor, Manus y Perplexity surgieron porque la capacidad de los modelos base se volvió de pronto suficientemente alta, y las empresas de producto la convirtieron, mediante un entorno o un contenedor, en un flujo de trabajo que el usuario puede percibir. Las empresas de aplicaciones no crean inteligencia de la nada: en el momento oportuno, recogen los beneficios del excedente de la investigación.

\begin{importantbox}{La tarea de las empresas de aplicaciones: construir contenedores de inteligencia}
Si el modelo base es el nivel del agua que sube, la tarea de las empresas de aplicaciones no es fingir que crearon el mar, sino construir bien el puerto, las tuberías y las llaves de agua. Cuanto más cerca esté el contenedor de las tareas reales, más clara sea la retroalimentación y más controlable el costo, más posibilidades tendrá de consolidar barreras de entrada.
\end{importantbox}

\subsection{El lugar de AI for Science y Robotics}

Esta subsección examina el tramo final de la hoja de ruta. Guangmi considera que en 2026/2027 podría darse una explosión de AI for Science, porque los problemas científicos tienen alto valor, retroalimentación verificable y un enorme espacio de mercado. Robotics está más cerca del mundo físico, pero él se mantiene escéptico ante el enfoque actual de los robotics foundation model: si solo se monta un VLM/VLA sobre un robot sin resolver de verdad los problemas de acción, datos, retroalimentación y cuerpo físico, el enfoque no llega a lo esencial.

\begin{warningbox}{No hay que mitificar Robotics antes de tiempo}
Los robots son un destino final importante en la ruta hacia la AGI, pero sus datos, su hardware, su retroalimentación física y sus restricciones de seguridad son mucho más complejos que los del mundo digital. Frente a obtener primero retroalimentación estable en los entornos de Coding y de Agent, apostar directamente por el mundo físico puede resultar más lento, más caro y más difícil de verificar.
\end{warningbox}

\subsection{Cortina de humo y línea principal: dónde ubicar la generación de imágenes a partir de texto}

Esta subsección aborda un juicio que se presta a controversia. Guangmi mencionó que la generación de imágenes a partir de texto podría ser una cortina de humo de OpenAI. No quiere decir que la generación de imágenes carezca de valor, sino que puede atraer demasiada atención sin representar necesariamente la máxima prioridad de la línea principal de la inteligencia. Para un fundamentalista de la AGI, la generación de imágenes, el video corto y los éxitos de consumo masivo deben volver a la misma pregunta: ¿hace que el modelo sea más inteligente, más capaz de actuar y más capaz de aprender?

\begin{knowledgebox}{Nota de clase: filtrar los temas de moda en IA según si aumentan la capacidad de acción}
Un tema de moda en IA puede ser muy comercial, muy llamativo y muy viral, pero si no profundiza la capacidad del modelo para entender el mundo, operar herramientas, obtener retroalimentación o transferir lo aprendido, no necesariamente pertenece a la línea principal de la AGI. El valor del informe trimestral está precisamente en usar la línea principal para filtrar los temas de moda.
\end{knowledgebox}

\subsection{Resumen de la sección}

El AGI roadmap de Guangmi sitúa a ChatGPT al pie de la montaña y coloca los entornos de acción y el descubrimiento científico en las cumbres principales que siguen. Esta hoja de ruta no necesariamente predice bien los plazos, pero ofrece un marco de juicio: cuanto más permita algo que el modelo actúe, verifique, aprenda y modifique el mundo, más cerca está de la línea principal.
